%% Generated by tools/gen_error_codes.py from docs/errors/ of the compiler;
%% edit the compiler's pages or tools/error_themes.txt, then rerun it.

\clearpage
\section{Control flow and pattern matching}
\label{sec:codes:flow}

The book's chapter \textit{Control flows} specifies the rules behind these codes.

\begin{longtable}{@{}l p{0.78\linewidth} r@{}}
\toprule Code & Message & Page\\ \midrule \endhead
\hyperref[sec:codes:e4006]{E4006} & array pattern match len \errhole{} never matches the value len \errhole{} & \pageref{sec:codes:e4006}\\
\hyperref[sec:codes:e4007]{E4007} & break statement is not within a loop & \pageref{sec:codes:e4007}\\
\hyperref[sec:codes:e4008]{E4008} & cannot break loop inside a scope guard & \pageref{sec:codes:e4008}\\
\hyperref[sec:codes:e4009]{E4009} & continue statement is not within a loop & \pageref{sec:codes:e4009}\\
\hyperref[sec:codes:e4010]{E4010} & cannot continue loop inside a scope guard & \pageref{sec:codes:e4010}\\
\hyperref[sec:codes:e4059]{E4059} & expected a name expression & \pageref{sec:codes:e4059}\\
\hyperref[sec:codes:e4067]{E4067} & decorators are forbidden for index iterator, on \errhole{} iteration & \pageref{sec:codes:e4067}\\
\hyperref[sec:codes:e4072]{E4072} & overriding the for loop operation on a class type \errhole{} & \pageref{sec:codes:e4072}\\
\hyperref[sec:codes:e4081]{E4081} & the if condition has no else value, so it must produce a void value not a \errhole{} & \pageref{sec:codes:e4081}\\
\hyperref[sec:codes:e4100]{E4100} & infinite loop with always true test is not allowed, rather use a loop construct & \pageref{sec:codes:e4100}\\
\hyperref[sec:codes:e4123]{E4123} & the pattern matching has no default case, so each branch must produce a void value not a \errhole{} & \pageref{sec:codes:e4123}\\
\hyperref[sec:codes:e4143]{E4143} & loop test is always false, loop is never entered & \pageref{sec:codes:e4143}\\
\hyperref[sec:codes:e4177]{E4177} & option matcher \errhole{} only applies to option values not \errhole{} & \pageref{sec:codes:e4177}\\
\hyperref[sec:codes:e4193]{E4193} & the pattern \errhole{} with value \errhole{} is refutable & \pageref{sec:codes:e4193}\\
\hyperref[sec:codes:e4197]{E4197} & return statement is not within a function & \pageref{sec:codes:e4197}\\
\hyperref[sec:codes:e4198]{E4198} & cannot return inside a scope guard & \pageref{sec:codes:e4198}\\
\hyperref[sec:codes:e4216]{E4216} & pattern \errhole{} only matches tuple values, not \errhole{} & \pageref{sec:codes:e4216}\\
\hyperref[sec:codes:e4239]{E4239} & undefined cte for loop operator with \errhole{} iterator for type \errhole{} & \pageref{sec:codes:e4239}\\
\hyperref[sec:codes:e4244]{E4244} & undefined for loop operator with \errhole{} iterator for type \errhole{} & \pageref{sec:codes:e4244}\\
\hyperref[sec:codes:e4259]{E4259} & matcher expression is never evaluated & \pageref{sec:codes:e4259}\\
\hyperref[sec:codes:e4260]{E4260} & unreachable statement & \pageref{sec:codes:e4260}\\
\hyperref[sec:codes:e4266]{E4266} & conditional test is always evaluated to false, branch is never entered & \pageref{sec:codes:e4266}\\
\hyperref[sec:codes:e4268]{E4268} & conditional test is always evaluated to true, branch is always entered & \pageref{sec:codes:e4268}\\
\hyperref[sec:codes:e4269]{E4269} & useless runtime assertion for a test that is always true & \pageref{sec:codes:e4269}\\
\hyperref[sec:codes:e4318]{E4318} & the pattern matching over \errhole{} has no default case, so it must cover each of its members & \pageref{sec:codes:e4318}\\
\bottomrule
\end{longtable}

\errorcode{E4006}{array pattern match len \errhole{} never matches the value len \errhole{}}
\label{sec:codes:e4006}

Raised during \textbf{validation}, from \texttt{ValidateErrorMessage::\allowbreak{}ARRAY\_\allowbreak{}PATTERN\_\allowbreak{}SIZE} (\textit{src/\allowbreak{}ymirc/\allowbreak{}semantic/\allowbreak{}validator/\allowbreak{}errors.yr}).

\subsubsection*{What it means}

An array pattern matches a fixed number of elements, and the value it is matched against has a different, statically known length, so the pattern can never match.

\subsubsection*{Example}

\textit{test\_\allowbreak{}resources/\allowbreak{}pattern\_\allowbreak{}match/\allowbreak{}test13.yr}:

%% check: skip
\begin{lstlisting}[style=coloredverbatimError]
fn main () {
    let x = [1, 2, 3];
    let [a = 1, b] = x;
}
\end{lstlisting}

\begin{lstlisting}[style=bashVerb, breaklines=true, escapechar=~]
Error[E4006] : array pattern match len 2us never matches the value len 3us
 --> test_resources/pattern_match/test13.yr:(4,9)
 4  ~\lstbox{\textSFxi{}}~     let [a = 1, b] = x;
    ~\lstbox{\textSFv{}}~         ^
    ~\lstbox{\textSFii{}}~~\lstbox{\textSFx{}}~~\lstbox{\textSFx{}}~~\lstbox{\textSFx{}}~~\lstbox{\textSFx{}}~~\lstbox{\textSFx{}}~~\lstbox{\textSFx{}}~~\lstbox{\textSFx{}}~
\end{lstlisting}

\subsubsection*{How to fix it}

Give the pattern the length of the value, or use a slice pattern if a variable number of elements was meant.

\errorcode{E4007}{break statement is not within a loop}
\label{sec:codes:e4007}

Raised during \textbf{validation}, from \texttt{ValidateErrorMessage::\allowbreak{}BREAK\_\allowbreak{}NO\_\allowbreak{}LOOP} (\textit{src/\allowbreak{}ymirc/\allowbreak{}semantic/\allowbreak{}validator/\allowbreak{}errors.yr}).

\subsubsection*{What it means}

A \tokennolst{break} was written outside any loop.

\subsubsection*{Example}

\textit{test\_\allowbreak{}resources/\allowbreak{}diagnostics/\allowbreak{}test18.yr}:

%% check: skip
\begin{lstlisting}[style=coloredverbatimError]
fn main () {
    break;
}
\end{lstlisting}

\begin{lstlisting}[style=bashVerb, breaklines=true, escapechar=~]
Error[E4007] : break statement is not within a loop
 --> test_resources/diagnostics/test18.yr:(2,5)
 2  ~\lstbox{\textSFxi{}}~     break;
    ~\lstbox{\textSFv{}}~     ^^^^^
    ~\lstbox{\textSFii{}}~~\lstbox{\textSFx{}}~~\lstbox{\textSFx{}}~~\lstbox{\textSFx{}}~~\lstbox{\textSFx{}}~~\lstbox{\textSFx{}}~~\lstbox{\textSFx{}}~~\lstbox{\textSFx{}}~
\end{lstlisting}

\subsubsection*{How to fix it}

Remove it, or move it inside the loop it was meant to leave. To leave a function early, use \tokennolst{return}.

\errorcode{E4008}{cannot break loop inside a scope guard}
\label{sec:codes:e4008}

Raised during \textbf{validation}, from \texttt{ValidateErrorMessage::\allowbreak{}BREAK\_\allowbreak{}SCOPE\_\allowbreak{}GUARD} (\textit{src/\allowbreak{}ymirc/\allowbreak{}semantic/\allowbreak{}validator/\allowbreak{}errors.yr}).

\subsubsection*{What it means}

A \tokennolst{break} inside a scope guard would leave the loop while the guard is running. A guard runs as part of leaving the scope it guards, so jumping out of it would abandon a scope exit half-way through.

\subsubsection*{Example}

\textit{test\_\allowbreak{}resources/\allowbreak{}diagnostics/\allowbreak{}test19.yr}:

%% check: skip
\begin{lstlisting}[style=coloredverbatimError]
fn main () {
    loop {
        {
        } exit {
            break;
        }
    }
}
\end{lstlisting}

\begin{lstlisting}[style=bashVerb, breaklines=true, escapechar=~]
Error[E4008] : cannot break loop inside a scope guard
 --> test_resources/diagnostics/test19.yr:(4,11)
 4  ~\lstbox{\textSFxi{}}~         } exit {
    ~\lstbox{\textSFv{}}~           ^^^^
   ...
 --> test_resources/diagnostics/test19.yr:(5,13)
 5  ~\lstbox{\textSFxi{}}~             break;
    ~\lstbox{\textSFv{}}~             ^^^^^
    ~\lstbox{\textSFii{}}~~\lstbox{\textSFx{}}~~\lstbox{\textSFx{}}~~\lstbox{\textSFx{}}~~\lstbox{\textSFx{}}~~\lstbox{\textSFx{}}~~\lstbox{\textSFx{}}~~\lstbox{\textSFx{}}~
\end{lstlisting}

\subsubsection*{How to fix it}

Move the \tokennolst{break} out of the guard body. A guard that has to influence the loop should set a flag the loop tests instead.

\errorcode{E4009}{continue statement is not within a loop}
\label{sec:codes:e4009}

Raised during \textbf{validation}, from \texttt{ValidateErrorMessage::\allowbreak{}CONTINUE\_\allowbreak{}NO\_\allowbreak{}LOOP} (\textit{src/\allowbreak{}ymirc/\allowbreak{}semantic/\allowbreak{}validator/\allowbreak{}errors.yr}).

\subsubsection*{What it means}

A \tokennolst{continue} was written outside any loop.

\subsubsection*{Example}

\textit{test\_\allowbreak{}resources/\allowbreak{}diagnostics/\allowbreak{}test20.yr}:

%% check: skip
\begin{lstlisting}[style=coloredverbatimError]
fn main () {
    continue;
}
\end{lstlisting}

\begin{lstlisting}[style=bashVerb, breaklines=true, escapechar=~]
Error[E4009] : continue statement is not within a loop
 --> test_resources/diagnostics/test20.yr:(2,5)
 2  ~\lstbox{\textSFxi{}}~     continue;
    ~\lstbox{\textSFv{}}~     ^^^^^^^^
    ~\lstbox{\textSFii{}}~~\lstbox{\textSFx{}}~~\lstbox{\textSFx{}}~~\lstbox{\textSFx{}}~~\lstbox{\textSFx{}}~~\lstbox{\textSFx{}}~~\lstbox{\textSFx{}}~~\lstbox{\textSFx{}}~
\end{lstlisting}

\subsubsection*{How to fix it}

Remove it, or move it inside the loop it was meant to restart.

\errorcode{E4010}{cannot continue loop inside a scope guard}
\label{sec:codes:e4010}

Raised during \textbf{validation}, from \texttt{ValidateErrorMessage::\allowbreak{}CONTINUE\_\allowbreak{}SCOPE\_\allowbreak{}GUARD} (\textit{src/\allowbreak{}ymirc/\allowbreak{}semantic/\allowbreak{}validator/\allowbreak{}errors.yr}).

\subsubsection*{What it means}

A \tokennolst{continue} inside a scope guard would jump out of the guard while it is running, cf. \hyperref[sec:codes:e4008]{E4008}.

\subsubsection*{Example}

\textit{test\_\allowbreak{}resources/\allowbreak{}diagnostics/\allowbreak{}test21.yr}:

%% check: skip
\begin{lstlisting}[style=coloredverbatimError]
fn main () {
    loop {
        {
        } exit {
            continue;
        }
    }
}
\end{lstlisting}

\begin{lstlisting}[style=bashVerb, breaklines=true, escapechar=~]
Error[E4010] : cannot continue loop inside a scope guard
 --> test_resources/diagnostics/test21.yr:(4,11)
 4  ~\lstbox{\textSFxi{}}~         } exit {
    ~\lstbox{\textSFv{}}~           ^^^^
   ...
 --> test_resources/diagnostics/test21.yr:(5,13)
 5  ~\lstbox{\textSFxi{}}~             continue;
    ~\lstbox{\textSFv{}}~             ^^^^^^^^
    ~\lstbox{\textSFii{}}~~\lstbox{\textSFx{}}~~\lstbox{\textSFx{}}~~\lstbox{\textSFx{}}~~\lstbox{\textSFx{}}~~\lstbox{\textSFx{}}~~\lstbox{\textSFx{}}~~\lstbox{\textSFx{}}~
\end{lstlisting}

\subsubsection*{How to fix it}

Move the \tokennolst{continue} out of the guard body.

\errorcode{E4059}{expected a name expression}
\label{sec:codes:e4059}

Raised during \textbf{validation}, from \texttt{ValidateErrorMessage::\allowbreak{}FIELD\_\allowbreak{}MATCHER\_\allowbreak{}NO\_\allowbreak{}NAME} (\textit{src/\allowbreak{}ymirc/\allowbreak{}semantic/\allowbreak{}validator/\allowbreak{}errors.yr}).

\subsubsection*{What it means}

A field pattern expected the name of a field and got something that is not a name.

\subsubsection*{Example}

\textit{test\_\allowbreak{}resources/\allowbreak{}diagnostics/\allowbreak{}test131.yr}:

%% check: skip
\begin{lstlisting}[style=coloredverbatimError]
// a record pattern with a field given by position
record Point {
    pub let x : i32;
    pub self (x : i32) with x = x {}
}

fn main () {
    let p = Point (1);
    match p {
        Point (1) => {}
        _ => {}
    }
}
\end{lstlisting}

\begin{lstlisting}[style=bashVerb, breaklines=true, escapechar=~]
Error[E4059] : expected a name expression
 --> test_resources/diagnostics/test131.yr:(10,16)
10  ~\lstbox{\textSFxi{}}~         Point (1) => {}
    ~\lstbox{\textSFv{}}~                ^
    ~\lstbox{\textSFxi{}}~
    = when validating pattern matching with value p of type test131::Point -- (test_resources/diagnostics/test131.yr:(9,5))
    = when validating test131::main -- (test_resources/diagnostics/test131.yr:(7,4))
    ~\lstbox{\textSFii{}}~~\lstbox{\textSFx{}}~~\lstbox{\textSFx{}}~~\lstbox{\textSFx{}}~~\lstbox{\textSFx{}}~~\lstbox{\textSFx{}}~~\lstbox{\textSFx{}}~~\lstbox{\textSFx{}}~
\end{lstlisting}

\subsubsection*{How to fix it}

Write the field name. A pattern matching by position rather than by name uses the tuple form instead.

\errorcode{E4067}{decorators are forbidden for index iterator, on \errhole{} iteration}
\label{sec:codes:e4067}

Raised during \textbf{validation}, from \texttt{ValidateErrorMessage::\allowbreak{}FORBID\_\allowbreak{}DECO\_\allowbreak{}FOR\_\allowbreak{}LOOP} (\textit{src/\allowbreak{}ymirc/\allowbreak{}semantic/\allowbreak{}validator/\allowbreak{}errors.yr}).

\subsubsection*{What it means}

A decorator was written on the index variable of a \tokennolst{for} loop. The index is produced by the loop itself, so its mutability is not the loop body's to choose.

\subsubsection*{Example}

\textit{test\_\allowbreak{}resources/\allowbreak{}for\_\allowbreak{}loops/\allowbreak{}arrays/\allowbreak{}test12.yr}:

%% check: skip
\begin{lstlisting}[style=coloredverbatimError]
fn main () {
    let a = [1, 2, 3];
    for mut i, _ in a {
        i;
    }
}
\end{lstlisting}

\begin{lstlisting}[style=bashVerb, breaklines=true, escapechar=~]
Error[E4067] : decorators are forbidden for index iterator, on [i32] iteration
 --> test_resources/for_loops/arrays/test12.yr:(3,9)
 3  ~\lstbox{\textSFxi{}}~     for mut i, _ in a {
    ~\lstbox{\textSFv{}}~         ^^^ ^
    ~\lstbox{\textSFii{}}~~\lstbox{\textSFx{}}~~\lstbox{\textSFx{}}~~\lstbox{\textSFx{}}~~\lstbox{\textSFx{}}~~\lstbox{\textSFx{}}~~\lstbox{\textSFx{}}~~\lstbox{\textSFx{}}~
\end{lstlisting}

\subsubsection*{How to fix it}

Remove the decorator. Copy the index into a local variable if a mutable one is needed.

\errorcode{E4072}{overriding the for loop operation on a class type \errhole{}}
\label{sec:codes:e4072}

Raised during \textbf{validation}, from \texttt{ValidateErrorMessage::\allowbreak{}FOR\_\allowbreak{}LOOP\_\allowbreak{}CPTR} (\textit{src/\allowbreak{}ymirc/\allowbreak{}semantic/\allowbreak{}validator/\allowbreak{}errors.yr}).

\subsubsection*{What it means}

The \tokennolst{for} loop operators are being overridden on a class type; this reports the context in which the override is validated.

\subsubsection*{How to fix it}

Check the \tokennolst{begin}, \tokennolst{next} and \tokennolst{get} methods the loop needs, which the notes under this diagnostic point at.

\errorcode{E4081}{the if condition has no else value, so it must produce a void value not a \errhole{}}
\label{sec:codes:e4081}

Raised during \textbf{validation}, from \texttt{ValidateErrorMessage::\allowbreak{}IF\_\allowbreak{}COND\_\allowbreak{}NOT\_\allowbreak{}COMPLETE} (\textit{src/\allowbreak{}ymirc/\allowbreak{}semantic/\allowbreak{}validator/\allowbreak{}errors.yr}).

\subsubsection*{What it means}

An \tokennolst{if} used as an expression produces a value on the branch that is taken, but has no \tokennolst{else}, so there would be nothing to produce when the condition is false.

\subsubsection*{Example}

\textit{test\_\allowbreak{}resources/\allowbreak{}control\_\allowbreak{}flow/\allowbreak{}test16.yr}:

%% check: skip
\begin{lstlisting}[style=coloredverbatimError]
fn main () {
    let mut cond = false;
    let a = if cond {
        12
    }; // error, no else value

    let b = if cond {
        12
    } else {
        "str"
    }; // error, incompatible i32 and [c8]

    let d = if cond {
        12
    } else {
        11u32
    }; // ok, implicit cast of 11u32 to i32
}
\end{lstlisting}

\begin{lstlisting}[style=bashVerb, breaklines=true, escapechar=~]
Error[E4081] : the if condition has no else value, so it must produce a void value not a i32
 --> test_resources/control_flow/test16.yr:(3,13)
 3  ~\lstbox{\textSFxi{}}~     let a = if cond {
    ~\lstbox{\textSFv{}}~             ^^
   ...
 --> test_resources/control_flow/test16.yr:(4,9)
 4  ~\lstbox{\textSFxi{}}~         12
    ~\lstbox{\textSFv{}}~         ^^
    ~\lstbox{\textSFii{}}~~\lstbox{\textSFx{}}~~\lstbox{\textSFx{}}~~\lstbox{\textSFx{}}~~\lstbox{\textSFx{}}~~\lstbox{\textSFx{}}~~\lstbox{\textSFx{}}~~\lstbox{\textSFx{}}~
\end{lstlisting}

\subsubsection*{How to fix it}

Add an \tokennolst{else} branch producing a value of the same type, or make the \tokennolst{if} a statement by ending the branch with something that produces nothing.

\errorcode{E4100}{infinite loop with always true test is not allowed, rather use a loop construct}
\label{sec:codes:e4100}

Raised during \textbf{validation}, from \texttt{ValidateErrorMessage::\allowbreak{}INFINITE\_\allowbreak{}LOOP} (\textit{src/\allowbreak{}ymirc/\allowbreak{}semantic/\allowbreak{}validator/\allowbreak{}errors.yr}).

\subsubsection*{What it means}

A \tokennolst{while} loop has a condition that is always true. Written that way, the condition suggests a test that does not exist.

\subsubsection*{Example}

\textit{test\_\allowbreak{}resources/\allowbreak{}control\_\allowbreak{}flow/\allowbreak{}test3.yr}:

%% check: skip
\begin{lstlisting}[style=coloredverbatimError]
fn main () {
    let a = [1, 2, 3];
    if let [x, y, z] = a { // error, len of /a/ is /cte/
        x;
        y;
        z;
    }

    if (a.len <= 3) {
    }

    if (a.len >= 4) {
    }
    
    while a.len <= 3 {
    }

    while a.len >= 4 {
    }

}
\end{lstlisting}

\begin{lstlisting}[style=bashVerb, breaklines=true, escapechar=~]
Error[E4100] : infinite loop with always true test is not allowed, rather use a loop construct
 --> test_resources/control_flow/test3.yr:(15,17)
15  ~\lstbox{\textSFxi{}}~     while a.len <= 3 {
    ~\lstbox{\textSFv{}}~                 ^^
    ~\lstbox{\textSFii{}}~~\lstbox{\textSFx{}}~~\lstbox{\textSFx{}}~~\lstbox{\textSFx{}}~~\lstbox{\textSFx{}}~~\lstbox{\textSFx{}}~~\lstbox{\textSFx{}}~~\lstbox{\textSFx{}}~
\end{lstlisting}

\subsubsection*{How to fix it}

Use \tokennolst{loop}, which says the loop is infinite on purpose.

\errorcode{E4123}{the pattern matching has no default case, so each branch must produce a void value not a \errhole{}}
\label{sec:codes:e4123}

Raised during \textbf{validation}, from \texttt{ValidateErrorMessage::\allowbreak{}MATCH\_\allowbreak{}NOT\_\allowbreak{}COMPLETE} (\textit{src/\allowbreak{}ymirc/\allowbreak{}semantic/\allowbreak{}validator/\allowbreak{}errors.yr}).

\subsubsection*{What it means}

A pattern matching used as an expression has no default case, so a value not covered by any pattern would have nothing to produce.

\subsubsection*{Example}

\textit{test\_\allowbreak{}resources/\allowbreak{}diagnostics/\allowbreak{}test52.yr}:

%% check: skip
\begin{lstlisting}[style=coloredverbatimError]
fn main () {
    let x = 1;
    let _ = match x {
        1 => { 2 }
    };
}
\end{lstlisting}

\begin{lstlisting}[style=bashVerb, breaklines=true, escapechar=~]
Error[E4123] : the pattern matching has no default case, so each branch must produce a void value not a i32
 --> test_resources/diagnostics/test52.yr:(4,14)
 4  ~\lstbox{\textSFxi{}}~         1 => { 2 }
    ~\lstbox{\textSFv{}}~              ^
    ~\lstbox{\textSFxi{}}~
    = when validating pattern matching with value x of type i32 -- (test_resources/diagnostics/test52.yr:(3,13))
    = when validating test52::main -- (test_resources/diagnostics/test52.yr:(1,4))
    ~\lstbox{\textSFii{}}~~\lstbox{\textSFx{}}~~\lstbox{\textSFx{}}~~\lstbox{\textSFx{}}~~\lstbox{\textSFx{}}~~\lstbox{\textSFx{}}~~\lstbox{\textSFx{}}~~\lstbox{\textSFx{}}~
\end{lstlisting}

\subsubsection*{How to fix it}

Add a default case producing a value of the same type, or make the match a statement by having its branches produce nothing.

A match over a sealed type -- a union or an enum -- is not reported here: it must cover the type whatever its arms produce, and is reported by \hyperref[sec:codes:e4318]{E4318} instead.

\errorcode{E4143}{loop test is always false, loop is never entered}
\label{sec:codes:e4143}

Raised during \textbf{validation}, from \texttt{ValidateErrorMessage::\allowbreak{}NEVER\_\allowbreak{}ENTERED\_\allowbreak{}LOOP} (\textit{src/\allowbreak{}ymirc/\allowbreak{}semantic/\allowbreak{}validator/\allowbreak{}errors.yr}).

\subsubsection*{What it means}

A loop condition is always false, so the body never runs.

\subsubsection*{Example}

\textit{test\_\allowbreak{}resources/\allowbreak{}control\_\allowbreak{}flow/\allowbreak{}test3.yr}:

%% check: skip
\begin{lstlisting}[style=coloredverbatimError]
fn main () {
    let a = [1, 2, 3];
    if let [x, y, z] = a { // error, len of /a/ is /cte/
        x;
        y;
        z;
    }

    if (a.len <= 3) {
    }

    if (a.len >= 4) {
    }
    
    while a.len <= 3 {
    }

    while a.len >= 4 {
    }

}
\end{lstlisting}

\begin{lstlisting}[style=bashVerb, breaklines=true, escapechar=~]
Error[E4143] : loop test is always false, loop is never entered
 --> test_resources/control_flow/test3.yr:(18,17)
18  ~\lstbox{\textSFxi{}}~     while a.len >= 4 {
    ~\lstbox{\textSFv{}}~                 ^^
    ~\lstbox{\textSFii{}}~~\lstbox{\textSFx{}}~~\lstbox{\textSFx{}}~~\lstbox{\textSFx{}}~~\lstbox{\textSFx{}}~~\lstbox{\textSFx{}}~~\lstbox{\textSFx{}}~~\lstbox{\textSFx{}}~
\end{lstlisting}

\subsubsection*{How to fix it}

Remove the loop, or fix the condition. This most often means the variable the condition tests is not the one the loop updates.

\errorcode{E4177}{option matcher \errhole{} only applies to option values not \errhole{}}
\label{sec:codes:e4177}

Raised during \textbf{validation}, from \texttt{ValidateErrorMessage::\allowbreak{}OPTION\_\allowbreak{}MATCHER} (\textit{src/\allowbreak{}ymirc/\allowbreak{}semantic/\allowbreak{}validator/\allowbreak{}errors.yr}).

\subsubsection*{What it means}

An option pattern was matched against something that is not an option.

\subsubsection*{Example}

\textit{test\_\allowbreak{}resources/\allowbreak{}diagnostics/\allowbreak{}test132.yr}:

%% check: skip
\begin{lstlisting}[style=coloredverbatimError]
// an option pattern matched against a value that is not an option
fn main () {
    match 1 {
        Ok (x) => { let _ = x; }
        _ => {}
    }
}
\end{lstlisting}

\begin{lstlisting}[style=bashVerb, breaklines=true, escapechar=~]
Error[E4177] : option matcher Ok --> test_resources/diagnostics/test132.yr:(4,9) only applies to option values not i32
 --> test_resources/diagnostics/test132.yr:(4,9)
 4  ~\lstbox{\textSFxi{}}~         Ok (x) => { let _ = x; }
    ~\lstbox{\textSFv{}}~         ^^
    ~\lstbox{\textSFxi{}}~
    = when validating pattern matching with value 1 of type i32 -- (test_resources/diagnostics/test132.yr:(3,5))
    = when validating test132::main -- (test_resources/diagnostics/test132.yr:(2,4))
    ~\lstbox{\textSFii{}}~~\lstbox{\textSFx{}}~~\lstbox{\textSFx{}}~~\lstbox{\textSFx{}}~~\lstbox{\textSFx{}}~~\lstbox{\textSFx{}}~~\lstbox{\textSFx{}}~~\lstbox{\textSFx{}}~
\end{lstlisting}

\subsubsection*{How to fix it}

Match the option itself, or use a pattern for the type at hand.

\errorcode{E4193}{the pattern \errhole{} with value \errhole{} is refutable}
\label{sec:codes:e4193}

Raised during \textbf{validation}, from \texttt{ValidateErrorMessage::\allowbreak{}PATTERN\_\allowbreak{}IS\_\allowbreak{}REFUTABLE} (\textit{src/\allowbreak{}ymirc/\allowbreak{}semantic/\allowbreak{}validator/\allowbreak{}errors.yr}).

\subsubsection*{What it means}

A pattern is used where every value has to match -- a destructuring binding, for instance -- and this one can fail.

\subsubsection*{Example}

\textit{test\_\allowbreak{}resources/\allowbreak{}control\_\allowbreak{}flow/\allowbreak{}test8.yr}:

%% check: skip
\begin{lstlisting}[style=coloredverbatimError]
fn foo ()-> i32?;

fn main () {
    let Ok (x : i32) = foo (); // error, refutable since it may fail if /foo/ returns /none/
}
\end{lstlisting}

\begin{lstlisting}[style=bashVerb, breaklines=true, escapechar=~]
Error[E4193] : the pattern Ok(x: i32) with value test8::foo () is refutable
 --> test_resources/control_flow/test8.yr:(4,5)
 4  ~\lstbox{\textSFxi{}}~     let Ok (x : i32) = foo (); // error, refutable since it may fail if /foo/ returns /none/
    ~\lstbox{\textSFv{}}~     ^^^
    ~\lstbox{\textSFii{}}~~\lstbox{\textSFx{}}~~\lstbox{\textSFx{}}~~\lstbox{\textSFx{}}~~\lstbox{\textSFx{}}~~\lstbox{\textSFx{}}~~\lstbox{\textSFx{}}~~\lstbox{\textSFx{}}~
\end{lstlisting}

\subsubsection*{How to fix it}

Use a pattern that always matches, or move the destructuring into a \tokennolst{match} with a default case.

\errorcode{E4197}{return statement is not within a function}
\label{sec:codes:e4197}

Raised during \textbf{validation}, from \texttt{ValidateErrorMessage::\allowbreak{}RETURN\_\allowbreak{}NO\_\allowbreak{}FUNCTION} (\textit{src/\allowbreak{}ymirc/\allowbreak{}semantic/\allowbreak{}validator/\allowbreak{}errors.yr}).

\subsubsection*{What it means}

A \tokennolst{return} was written outside any function.

\subsubsection*{Example}

\textit{test\_\allowbreak{}resources/\allowbreak{}diagnostics/\allowbreak{}test69.yr}:

%% check: skip
\begin{lstlisting}[style=coloredverbatimError]
lazy X : i32 = { return 1; };

fn main () { X; }
\end{lstlisting}

\begin{lstlisting}[style=bashVerb, breaklines=true, escapechar=~]
Error[E4197] : return statement is not within a function
 --> test_resources/diagnostics/test69.yr:(1,18)
 1  ~\lstbox{\textSFxi{}}~ lazy X : i32 = { return 1; };
    ~\lstbox{\textSFv{}}~                  ^^^^^^
    ~\lstbox{\textSFxi{}}~
    = when validating -- (test_resources/diagnostics/test69.yr:(1,6))
    ~\lstbox{\textSFii{}}~~\lstbox{\textSFx{}}~~\lstbox{\textSFx{}}~~\lstbox{\textSFx{}}~~\lstbox{\textSFx{}}~~\lstbox{\textSFx{}}~~\lstbox{\textSFx{}}~~\lstbox{\textSFx{}}~
Error[E4104] : symbol test69::X has errors
 --> test_resources/diagnostics/test69.yr:(3,14)
 1  ~\lstbox{\textSFxi{}}~ lazy X : i32 = { return 1; };
    ~\lstbox{\textSFv{}}~                  ------  error: return statement is not within a function
 2  ~\lstbox{\textSFxi{}}~
 3  ~\lstbox{\textSFxi{}}~ fn main () { X; }
    ~\lstbox{\textSFv{}}~              ^
    ~\lstbox{\textSFxi{}}~
    = when validating -- (test_resources/diagnostics/test69.yr:(1,6))
    = when validating test69::main -- (test_resources/diagnostics/test69.yr:(3,4))
    ~\lstbox{\textSFii{}}~~\lstbox{\textSFx{}}~~\lstbox{\textSFx{}}~~\lstbox{\textSFx{}}~~\lstbox{\textSFx{}}~~\lstbox{\textSFx{}}~~\lstbox{\textSFx{}}~~\lstbox{\textSFx{}}~
\end{lstlisting}

\subsubsection*{How to fix it}

Remove it, or move it into the function it was meant to leave.

\errorcode{E4198}{cannot return inside a scope guard}
\label{sec:codes:e4198}

Raised during \textbf{validation}, from \texttt{ValidateErrorMessage::\allowbreak{}RETURN\_\allowbreak{}SCOPE\_\allowbreak{}GUARD} (\textit{src/\allowbreak{}ymirc/\allowbreak{}semantic/\allowbreak{}validator/\allowbreak{}errors.yr}).

\subsubsection*{What it means}

A \tokennolst{return} inside a scope guard would leave the function while the guard is running, abandoning a scope exit half-way through, cf. \hyperref[sec:codes:e4008]{E4008}.

\subsubsection*{Example}

\textit{test\_\allowbreak{}resources/\allowbreak{}scope\_\allowbreak{}guards/\allowbreak{}test7.yr}:

%% check: skip
\begin{lstlisting}[style=coloredverbatimError]
fn main () {
    {

    } exit {
        {
            throw copy AssertError ("");
        } catch {
            _ => return ;
        }
    }
}
\end{lstlisting}

\begin{lstlisting}[style=bashVerb, breaklines=true, escapechar=~]
Error[E4198] : cannot return inside a scope guard
 --> test_resources/scope_guards/test7.yr:(4,7)
 4  ~\lstbox{\textSFxi{}}~     } exit {
    ~\lstbox{\textSFv{}}~       ^^^^
 --> test_resources/scope_guards/test7.yr:(8,18)
 8  ~\lstbox{\textSFxi{}}~             _ => return ;
    ~\lstbox{\textSFv{}}~                  ------
    ~\lstbox{\textSFxi{}}~
    = when validating catching pattern with type &(core::exception::assertion::AssertError) -- (test_resources/scope_guards/test7.yr:(7,11))
    = when validating test7::main -- (test_resources/scope_guards/test7.yr:(1,4))
    ~\lstbox{\textSFii{}}~~\lstbox{\textSFx{}}~~\lstbox{\textSFx{}}~~\lstbox{\textSFx{}}~~\lstbox{\textSFx{}}~~\lstbox{\textSFx{}}~~\lstbox{\textSFx{}}~~\lstbox{\textSFx{}}~
\end{lstlisting}

\subsubsection*{How to fix it}

Move the \tokennolst{return} out of the guard body.

\errorcode{E4216}{pattern \errhole{} only matches tuple values, not \errhole{}}
\label{sec:codes:e4216}

Raised during \textbf{validation}, from \texttt{ValidateErrorMessage::\allowbreak{}TUPLE\_\allowbreak{}PATTERN} (\textit{src/\allowbreak{}ymirc/\allowbreak{}semantic/\allowbreak{}validator/\allowbreak{}errors.yr}).

\subsubsection*{What it means}

A pattern that destructures a tuple was applied to a value that is not one. Both the pattern and the type it was matched against are in the message.

\subsubsection*{Example}

\textit{test\_\allowbreak{}resources/\allowbreak{}diagnostics/\allowbreak{}test133.yr}:

%% check: skip
\begin{lstlisting}[style=coloredverbatimError]
// a tuple pattern matched against a value that is not a tuple
fn main () {
    match 1 {
        (x, y) => { let _ = (x, y); }
        _ => {}
    }
}
\end{lstlisting}

\begin{lstlisting}[style=bashVerb, breaklines=true, escapechar=~]
Error[E4216] : pattern (x, y) only matches tuple values, not i32
 --> test_resources/diagnostics/test133.yr:(4,9)
 4  ~\lstbox{\textSFxi{}}~         (x, y) => { let _ = (x, y); }
    ~\lstbox{\textSFv{}}~         ^
    ~\lstbox{\textSFxi{}}~
    = when validating pattern matching with value 1 of type i32 -- (test_resources/diagnostics/test133.yr:(3,5))
    = when validating test133::main -- (test_resources/diagnostics/test133.yr:(2,4))
    ~\lstbox{\textSFii{}}~~\lstbox{\textSFx{}}~~\lstbox{\textSFx{}}~~\lstbox{\textSFx{}}~~\lstbox{\textSFx{}}~~\lstbox{\textSFx{}}~~\lstbox{\textSFx{}}~~\lstbox{\textSFx{}}~
\end{lstlisting}

\subsubsection*{How to fix it}

Match a tuple, or use a pattern defined for the type at hand. A record is not a tuple: its fields are matched by name, not by position.

\errorcode{E4239}{undefined cte for loop operator with \errhole{} iterator for type \errhole{}}
\label{sec:codes:e4239}

Raised during \textbf{validation}, from \texttt{ValidateErrorMessage::\allowbreak{}UNDEF\_\allowbreak{}CTE\_\allowbreak{}FOR\_\allowbreak{}LOOP\_\allowbreak{}OPERATOR} (\textit{src/\allowbreak{}ymirc/\allowbreak{}semantic/\allowbreak{}validator/\allowbreak{}errors.yr}).

\subsubsection*{What it means}

A \tokennolst{cte for} loop was written over a type that defines no compile time iteration.

\subsubsection*{Example}

\textit{test\_\allowbreak{}resources/\allowbreak{}diagnostics/\allowbreak{}test123.yr}:

%% check: skip
\begin{lstlisting}[style=coloredverbatimError]
// a compile time loop over a slice, whose length is not known at compile time
fn main () {
    cte for i in copy [1, 2] {
        let _ = i;
    }
}
\end{lstlisting}

\begin{lstlisting}[style=bashVerb, breaklines=true, escapechar=~]
Error[E4239] : undefined cte for loop operator with 1 iterator for type mut [mut i32]
 --> test_resources/diagnostics/test123.yr:(3,9)
 3  ~\lstbox{\textSFxi{}}~     cte for i in copy [1, 2] {
    ~\lstbox{\textSFv{}}~         ^^^
    ~\lstbox{\textSFxi{}}~
    = when validating test123::main -- (test_resources/diagnostics/test123.yr:(2,4))
    ~\lstbox{\textSFii{}}~~\lstbox{\textSFx{}}~~\lstbox{\textSFx{}}~~\lstbox{\textSFx{}}~~\lstbox{\textSFx{}}~~\lstbox{\textSFx{}}~~\lstbox{\textSFx{}}~~\lstbox{\textSFx{}}~
\end{lstlisting}

\subsubsection*{How to fix it}

Iterate over a type that does -- a tuple, an array, or a compile time range.

\errorcode{E4244}{undefined for loop operator with \errhole{} iterator for type \errhole{}}
\label{sec:codes:e4244}

Raised during \textbf{validation}, from \texttt{ValidateErrorMessage::\allowbreak{}UNDEF\_\allowbreak{}FOR\_\allowbreak{}LOOP\_\allowbreak{}OPERATOR} (\textit{src/\allowbreak{}ymirc/\allowbreak{}semantic/\allowbreak{}validator/\allowbreak{}errors.yr}).

\subsubsection*{What it means}

A \tokennolst{for} loop was written over a type that defines no iteration. A type is iterable when it provides the \tokennolst{begin}, \tokennolst{next} and \tokennolst{get} operators.

\subsubsection*{Example}

\textit{test\_\allowbreak{}resources/\allowbreak{}diagnostics/\allowbreak{}test77.yr}:

%% check: skip
\begin{lstlisting}[style=coloredverbatimError]
fn main () {
    for i in 1 { i; }
}
\end{lstlisting}

\begin{lstlisting}[style=bashVerb, breaklines=true, escapechar=~]
Error[E4244] : undefined for loop operator with 1 iterator for type i32
 --> test_resources/diagnostics/test77.yr:(2,5)
 2  ~\lstbox{\textSFxi{}}~     for i in 1 { i; }
    ~\lstbox{\textSFv{}}~     ^^^
    ~\lstbox{\textSFii{}}~~\lstbox{\textSFx{}}~~\lstbox{\textSFx{}}~~\lstbox{\textSFx{}}~~\lstbox{\textSFx{}}~~\lstbox{\textSFx{}}~~\lstbox{\textSFx{}}~~\lstbox{\textSFx{}}~
\end{lstlisting}

\subsubsection*{How to fix it}

Define those operators on the type, or iterate over something that already has them.

\errorcode{E4259}{matcher expression is never evaluated}
\label{sec:codes:e4259}

Raised during \textbf{validation}, from \texttt{ValidateErrorMessage::\allowbreak{}UNREACHABLE\_\allowbreak{}MATCHER} (\textit{src/\allowbreak{}ymirc/\allowbreak{}semantic/\allowbreak{}validator/\allowbreak{}errors.yr}).

\subsubsection*{What it means}

A pattern in a \tokennolst{match} is covered by an earlier one, so it can never run. A default case raises it too when the arms before it already cover every member of the union, or every field of the enum, being matched.

\subsubsection*{How to fix it}

Remove it, or move it before the pattern that shadows it.

\errorcode{E4260}{unreachable statement}
\label{sec:codes:e4260}

Raised during \textbf{validation}, from \texttt{ValidateErrorMessage::\allowbreak{}UNREACHBLE\_\allowbreak{}STATEMENT} (\textit{src/\allowbreak{}ymirc/\allowbreak{}semantic/\allowbreak{}validator/\allowbreak{}errors.yr}).

\subsubsection*{What it means}

A statement follows one that never completes -- a \tokennolst{return}, a \tokennolst{break}, a \tokennolst{throw} -- so it can never run.

\subsubsection*{Example}

\textit{test\_\allowbreak{}resources/\allowbreak{}lit\_\allowbreak{}options/\allowbreak{}test16.yr}:

%% check: skip
\begin{lstlisting}[style=coloredverbatimError]
fn main () {
    loop {
        let _a_ = ({
            break 12;
        })?;
    };
}
\end{lstlisting}

\begin{lstlisting}[style=bashVerb, breaklines=true, escapechar=~]
Error[E4260] : unreachable statement
 --> test_resources/lit_options/test16.yr:(5,11)
 5  ~\lstbox{\textSFxi{}}~         })?;
    ~\lstbox{\textSFv{}}~           ^
    ~\lstbox{\textSFxi{}}~ Note : exiting scope here
    ~\lstbox{\textSFxi{}}~  --> test_resources/lit_options/test16.yr:(4,13)
    ~\lstbox{\textSFxi{}}~  4  ~\lstbox{\textSFxi{}}~             break 12;
    ~\lstbox{\textSFxi{}}~     ~\lstbox{\textSFv{}}~             ^^^^^
    ~\lstbox{\textSFxi{}}~     ~\lstbox{\textSFii{}}~~\lstbox{\textSFx{}}~~\lstbox{\textSFx{}}~~\lstbox{\textSFx{}}~~\lstbox{\textSFx{}}~~\lstbox{\textSFx{}}~~\lstbox{\textSFx{}}~~\lstbox{\textSFx{}}~
    ~\lstbox{\textSFii{}}~~\lstbox{\textSFx{}}~~\lstbox{\textSFx{}}~~\lstbox{\textSFx{}}~~\lstbox{\textSFx{}}~~\lstbox{\textSFx{}}~~\lstbox{\textSFvii{}}~~\lstbox{\textSFx{}}~
\end{lstlisting}

\subsubsection*{How to fix it}

Remove the statement, or move it before the one that leaves the block.

\errorcode{E4266}{conditional test is always evaluated to false, branch is never entered}
\label{sec:codes:e4266}

Raised during \textbf{validation}, from \texttt{ValidateErrorMessage::\allowbreak{}USELESS\_\allowbreak{}COND\_\allowbreak{}FALSE} (\textit{src/\allowbreak{}ymirc/\allowbreak{}semantic/\allowbreak{}validator/\allowbreak{}errors.yr}).

\subsubsection*{What it means}

A condition the compiler could evaluate is always false, so the branch it guards is dead.

\subsubsection*{Example}

\textit{test\_\allowbreak{}resources/\allowbreak{}control\_\allowbreak{}flow/\allowbreak{}test15.yr}:

%% check: skip
\begin{lstlisting}[style=coloredverbatimError]
fn main () {
    let a = [1, 2, 3];
    if (a.len <= 3) { // error, len of /a/ is /cte/
        println (a [2]);
    }

    if (a.len >= 4) {
        println (a [3]); // error, dead code
    }
}
\end{lstlisting}

\begin{lstlisting}[style=bashVerb, breaklines=true, escapechar=~]
Error[E4266] : conditional test is always evaluated to false, branch is never entered
 --> test_resources/control_flow/test15.yr:(7,15)
 7  ~\lstbox{\textSFxi{}}~     if (a.len >= 4) {
    ~\lstbox{\textSFv{}}~               ^^
    ~\lstbox{\textSFii{}}~~\lstbox{\textSFx{}}~~\lstbox{\textSFx{}}~~\lstbox{\textSFx{}}~~\lstbox{\textSFx{}}~~\lstbox{\textSFx{}}~~\lstbox{\textSFx{}}~~\lstbox{\textSFx{}}~
\end{lstlisting}

\subsubsection*{How to fix it}

Remove the branch, or fix the condition.

\errorcode{E4268}{conditional test is always evaluated to true, branch is always entered}
\label{sec:codes:e4268}

Raised during \textbf{validation}, from \texttt{ValidateErrorMessage::\allowbreak{}USELESS\_\allowbreak{}COND\_\allowbreak{}TRUE} (\textit{src/\allowbreak{}ymirc/\allowbreak{}semantic/\allowbreak{}validator/\allowbreak{}errors.yr}).

\subsubsection*{What it means}

A condition the compiler could evaluate is always true, so the branch it guards always runs and the test says nothing.

\subsubsection*{Example}

\textit{test\_\allowbreak{}resources/\allowbreak{}control\_\allowbreak{}flow/\allowbreak{}test15.yr}:

%% check: skip
\begin{lstlisting}[style=coloredverbatimError]
fn main () {
    let a = [1, 2, 3];
    if (a.len <= 3) { // error, len of /a/ is /cte/
        println (a [2]);
    }

    if (a.len >= 4) {
        println (a [3]); // error, dead code
    }
}
\end{lstlisting}

\begin{lstlisting}[style=bashVerb, breaklines=true, escapechar=~]
Error[E4268] : conditional test is always evaluated to true, branch is always entered
 --> test_resources/control_flow/test15.yr:(3,15)
 3  ~\lstbox{\textSFxi{}}~     if (a.len <= 3) { // error, len of /a/ is /cte/
    ~\lstbox{\textSFv{}}~               ^^
    ~\lstbox{\textSFii{}}~~\lstbox{\textSFx{}}~~\lstbox{\textSFx{}}~~\lstbox{\textSFx{}}~~\lstbox{\textSFx{}}~~\lstbox{\textSFx{}}~~\lstbox{\textSFx{}}~~\lstbox{\textSFx{}}~
\end{lstlisting}

\subsubsection*{How to fix it}

Remove the condition, or fix it.

\errorcode{E4269}{useless runtime assertion for a test that is always true}
\label{sec:codes:e4269}

Raised during \textbf{validation}, from \texttt{ValidateErrorMessage::\allowbreak{}USELESS\_\allowbreak{}RUNTIME\_\allowbreak{}ASSERT} (\textit{src/\allowbreak{}ymirc/\allowbreak{}semantic/\allowbreak{}validator/\allowbreak{}errors.yr}).

\subsubsection*{What it means}

A run time assertion tests something the compiler already proved true, so it can never fire.

\subsubsection*{Example}

\textit{test\_\allowbreak{}resources/\allowbreak{}diagnostics/\allowbreak{}test81.yr}:

%% check: skip
\begin{lstlisting}[style=coloredverbatimError]
fn main () {
    assert (true);
}
\end{lstlisting}

\begin{lstlisting}[style=bashVerb, breaklines=true, escapechar=~]
Error[E4269] : useless runtime assertion for a test that is always true
 --> test_resources/diagnostics/test81.yr:(2,5)
 2  ~\lstbox{\textSFxi{}}~     assert (true);
    ~\lstbox{\textSFv{}}~     ^^^^^^  ^^^^
    ~\lstbox{\textSFii{}}~~\lstbox{\textSFx{}}~~\lstbox{\textSFx{}}~~\lstbox{\textSFx{}}~~\lstbox{\textSFx{}}~~\lstbox{\textSFx{}}~~\lstbox{\textSFx{}}~~\lstbox{\textSFx{}}~
\end{lstlisting}

\subsubsection*{How to fix it}

Remove the assertion, or make it a \tokennolst{cte assert} if the compile time check is what was wanted.

\errorcode{E4318}{the pattern matching over \errhole{} has no default case, so it must cover each of its members}
\label{sec:codes:e4318}

Raised during \textbf{validation}, from \texttt{ValidateErrorMessage::\allowbreak{}MATCH\_\allowbreak{}SEALED\_\allowbreak{}NOT\_\allowbreak{}COMPLETE} (\textit{src/\allowbreak{}ymirc/\allowbreak{}semantic/\allowbreak{}validator/\allowbreak{}errors.yr}).

\subsubsection*{What it means}

A \tokennolst{match} over a union or an enum has no default case and leaves out a member of the union, or a field of the enum. A union and an enum list everything they can hold, so such a match must cover that list, whether its arms produce a value or not: adding a member to the type is then an error at every match that forgets it, instead of a value that silently falls through. The notes name each member or field no arm covers.

An arm carrying a guard, or a pattern filtering the fields of the member, decides on more than the member it names, so it does not count as covering it.

\subsubsection*{Example}

\textit{test\_\allowbreak{}resources/\allowbreak{}type\_\allowbreak{}union/\allowbreak{}test83.yr}:

%% check: skip
\begin{lstlisting}[style=coloredverbatimError]
// A match over a union must be complete even when its arms produce nothing, so
// the member no arm names is reported.

record X {
    pub let a : i32;

    pub self (a : i32)
        with a = a
    {}
}

record Y {
    pub let b : i32;

    pub self (b : i32)
        with b = b
    {}
}

union U
| X
| Y;

fn run (u : U) {
    match u {
        _ : X => {}
    }
}

fn main () {
    run (X (1));
}
\end{lstlisting}

\begin{lstlisting}[style=bashVerb, breaklines=true, escapechar=~]
Error[E4318] : the pattern matching over test83::U has no default case, so it must cover each of its members
 --> test_resources/type_union/test83.yr:(25,5)
25  ~\lstbox{\textSFxi{}}~     match u {
    ~\lstbox{\textSFv{}}~     ^^^^^
    ~\lstbox{\textSFxi{}}~
    = when validating pattern matching with value u of type test83::U -- (test_resources/type_union/test83.yr:(25,5))
    = when validating test83::run -- (test_resources/type_union/test83.yr:(24,4))
    ~\lstbox{\textSFii{}}~~\lstbox{\textSFx{}}~~\lstbox{\textSFx{}}~~\lstbox{\textSFx{}}~~\lstbox{\textSFx{}}~~\lstbox{\textSFx{}}~~\lstbox{\textSFx{}}~~\lstbox{\textSFx{}}~
Note : the member test83::Y is never matched
 --> test_resources/type_union/test83.yr:(25,5)
25  ~\lstbox{\textSFxi{}}~     match u {
    ~\lstbox{\textSFv{}}~     ^^^^^
    ~\lstbox{\textSFxi{}}~
    = when validating pattern matching with value u of type test83::U -- (test_resources/type_union/test83.yr:(25,5))
    = when validating test83::run -- (test_resources/type_union/test83.yr:(24,4))
    ~\lstbox{\textSFii{}}~~\lstbox{\textSFx{}}~~\lstbox{\textSFx{}}~~\lstbox{\textSFx{}}~~\lstbox{\textSFx{}}~~\lstbox{\textSFx{}}~~\lstbox{\textSFx{}}~~\lstbox{\textSFx{}}~
\end{lstlisting}

\subsubsection*{How to fix it}

Add an arm for each member or field the notes name, or a default case (\tokennolst{\_ =\textgreater{}}) when leaving them out is deliberate.
